%% =============================================================================
%% Jiuzhang Suanshu (Nine Chapters on the Mathematical Art) -- Fascicles 1-3
%% With Liu Hui's Commentary (263 CE)
%% EXPANDED EDITION -- Complete Chinese Text with Extensive Liu Hui Commentary
%% =============================================================================
\documentclass[11pt,a4paper]{article}

%% -------- Core Packages --------
\usepackage{fontspec}
\usepackage{polyglossia}
\usepackage{amsmath,amssymb,amsthm}
\usepackage{geometry}
\usepackage{graphicx}
\usepackage{xcolor}
\usepackage{titlesec}
\usepackage{fancyhdr}
\usepackage{enumitem}
\usepackage{array,booktabs,longtable,multirow}
\usepackage[draft]{hyperref}
\usepackage{tocloft}
\usepackage{indentfirst}
\usepackage{xeCJK}
\usepackage{verse}
\usepackage{lettrine}
\usepackage{microtype}
% \usepackage{quoting}  % disabled: package not available

%% -------- Page Geometry --------
\geometry{
  paperwidth=210mm,paperheight=297mm,
  left=24mm,right=24mm,top=28mm,bottom=28mm,
  headheight=14pt,footskip=30pt
}

%% -------- Language Setup --------
\setmainlanguage{english}
\setotherlanguage{french}

%% -------- Font Configuration --------
\setmainfont{Linux Libertine O}[
  Extension=.otf,
  UprightFont=*-Regular,
  BoldFont=*-Bold,
  ItalicFont=*-Italic,
  BoldItalicFont=*-BoldItalic,
  Renderer=HarfbuzzX  % Changed: Harfbuzz triggers CTAN fetches
]
\setsansfont{Linux Biolinum O}[
  Extension=.otf,
  UprightFont=*-Regular,
  BoldFont=*-Bold,
  ItalicFont=*-Italic,
  Renderer=HarfbuzzX  % Changed: Harfbuzz triggers CTAN fetches
]
\setmonofont{DejaVu Sans Mono}[Scale=MatchLowercase]

%% CJK Font Setup
\setCJKmainfont{NotoSansCJKsc-Regular.otf}[
  Path=fonts/noto-sc-extracted/,
  BoldFont=NotoSansCJKsc-Bold.otf
]
\setCJKsansfont{NotoSansCJKsc-Regular.otf}[
  Path=fonts/noto-sc-extracted/,
  BoldFont=NotoSansCJKsc-Bold.otf
]
\setCJKmonofont{NotoSansCJKsc-Regular.otf}[
  Path=fonts/noto-sc-extracted/,
  BoldFont=NotoSansCJKsc-Bold.otf
]

%% -------- Colors --------
\definecolor{titlecolor}{RGB}{45,55,72}
\definecolor{sectioncolor}{RGB}{55,65,81}
\definecolor{notecolor}{RGB}{120,53,15}
\definecolor{rulecolor}{RGB}{180,160,140}
\definecolor{commentarycolor}{RGB}{70,90,110}
\definecolor{problemcolor}{RGB}{30,60,90}
\definecolor{answercolor}{RGB}{40,80,60}
\definecolor{methodcolor}{RGB}{80,50,30}
\definecolor{prefacecolor}{RGB}{60,50,70}

%% -------- Section Formatting --------
\titleformat{\section}
  {\normalfont\Large\bfseries\color{sectioncolor}}
  {\thesection}{1em}{}
\titleformat{\subsection}
  {\normalfont\large\bfseries\color{sectioncolor}}
  {\thesubsection}{1em}{}
\titleformat{\subsubsection}
  {\normalfont\normalsize\bfseries\color{sectioncolor}}
  {\thesubsubsection}{1em}{}

%% -------- Header/Footer --------
\pagestyle{fancy}
\fancyhf{}
\fancyhead[L]{\small\textit{九章算術 -- 卷第一至卷第三}}
\fancyhead[R]{\small\thepage}
\renewcommand{\headrulewidth}{0.4pt}
\renewcommand{\footrulewidth}{0pt}

%% -------- Custom Environments --------
\newenvironment{problem}[1][]{%
  \par\vspace{0.6em}\noindent\textcolor{problemcolor}{\textbf{問\ignorespaces#1\unskip}\hspace{0.5em}}%
  \ignorespaces
}{\par\vspace{0.4em}}

\newenvironment{answer}{%
  \par\noindent\textcolor{answercolor}{\textbf{荅\,曰：}}%
  \ignorespaces
}{\par\vspace{0.3em}}

\newenvironment{method}{%
  \par\noindent\textcolor{methodcolor}{\textbf{術\,曰：}}%
  \ignorespaces
}{\par\vspace{0.4em}}

\newenvironment{commentary}{%
  \par\noindent\textcolor{commentarycolor}{\small\textit{（注）}}%
  \ignorespaces
}{\par\vspace{0.4em}}

\newenvironment{liuhuicomm}{%
  \par\vspace{0.3em}\noindent\textcolor{commentarycolor}{\small\textit{劉徽注：}}}%
  {\par\vspace{0.5em}}

\newenvironment{lichunfeng}{%
  \par\vspace{0.3em}\noindent\textcolor{commentarycolor}{\small\textit{李淳風等注：}}}%
  {\par\vspace{0.5em}}

\newenvironment{prefacetext}{%
  \par\vspace{0.5em}\noindent\textcolor{prefacecolor}{\small}%
  \ignorespaces
}{\par\vspace{0.5em}}

\newcommand{\liuhui}[1]{\textcolor{commentarycolor}{\small #1}}
\newcommand{\problemnum}[1]{\textbf{[\ignorespaces#1\unskip]}}
\newcommand{\zishu}[1]{\textit{\textcolor{notecolor}{#1}}}

%% -------- Hyperref Setup --------
\hypersetup{
  colorlinks=true,
  linkcolor=sectioncolor,
  citecolor=sectioncolor,
  urlcolor=sectioncolor,
  pdfencoding=auto,
  pdftitle={Jiuzhang Suanshu -- Volumes 1-3 with Liu Hui Commentary (Expanded)},
  pdfauthor={Anonymous; Commentary by Liu Hui},
}

%% -------- TOC Depth --------
\setcounter{tocdepth}{2}
\setcounter{secnumdepth}{2}

%% -------- Vertical Spacing --------
\setlength{\parskip}{0.3em}
\setlength{\parindent}{0pt}

%% =============================================================================
\begin{document}

%% -------- Title Page --------
\begin{titlepage}
\begin{center}
\vspace*{2.5cm}

{\Huge\bfseries\color{titlecolor} Ji\u{u}zh\={a}ng Su\`{a}nsh\`{u}}\\[0.6cm]
{\LARGE\bfseries\color{sectioncolor} 九章算術}\\[1.2cm]
{\Large\textit{Nine Chapters on the Mathematical Art}}\\[1.8cm]
{\large With Extensive Commentary by \textbf{Liu Hui} (劉徽, 263 CE)}\\[0.4cm]
{\large and Later Annotations by Li Chunfeng (李淳風, 656 CE)}\\[2.5cm]

\rule{0.6\textwidth}{0.8pt}\\[1.2cm]

{\Large\bfseries 卷第一至卷第三}\\[0.8cm]
{\large Volumes 1--3 (Expanded Edition)}\\[1.5cm]

\begin{tabular}{cl}
\textbf{卷一} & 方田 \quad Fangti\'{a}n \quad (Rectangular Fields)\\
\textbf{卷二} & 粟米 \quad S\`{u}m\u{i} \quad (Millet and Rice)\\
\textbf{卷三} & 衰分 \quad Cu\={i}f\={e}n \quad (Distributing by Proportion)\\
\end{tabular}

\vfill

{\small The most important classical Chinese mathematical text}\\[0.3cm]
{\small Compiled c.~200 BCE--200 CE \quad|\quad Commentary 263 CE}\\[0.5cm]
{\small\textit{Expanded edition with complete problems, methods, and full Liu Hui commentary}}\\[1cm]

\rule{0.4\textwidth}{0.4pt}

\end{center}
\end{titlepage}

%% -------- Table of Contents --------
\tableofcontents
\newpage


%% =============================================================================
\section*{劉徽九章算術注原序}
\addcontentsline{toc}{section}{原序：劉徽九章算術注}

\begin{center}
\rule{0.5\textwidth}{0.4pt}
\end{center}

\begin{prefacetext}
昔在包犧氏始畫八卦，以通神明之德，以類萬物之情，作九九之術以合六爻之變。
暨於黃帝神而化之，引而伸之，於是建曆紀，協律呂，用稽道原，然後兩儀四象精微之氣可得而效焉。
記稱隸首作數，其詳未之聞也。
按周公制禮而有九數，九數之流，則九章是矣。

往者暴秦焚書，經術散壞。自時厥後，漢北平侯張蒼、大司農中丞耿壽昌皆以善算命世。
蒼等因舊文之遺殘，各稱刪補。故校其目則與古或異，而所論者多近語也。

徽幼習九章，長再詳覽。觀陰陽之割裂，總算術之根源，探賾之暇，遂悟其意。
是以敢竭頑魯，采其所見，為之作注。
事類相推，各有攸歸，故枝條雖分而同本榦者，知發其一端而已。
又所析理以辭，解體用圖，庶亦約而能周，通而不黷，覽之者思過半矣。
且算在六藝，古者以賓興賢能，教習國子。雖曰九數，其能窮纖入微，探測無方。
至於以法相傳，亦猶規矩度量可得而共，非特難為也。
當今好之者寡，故世雖多通才達學，而未必能綜於此耳。

周官大司徒職，夏至日中立八尺之表，其景尺有五寸，謂之地中。
說云，南戴日下萬五千里。夫云爾者，以術推之。
按九章立四表望遠及因木望山之術，皆端旁互見，無有超邈若斯之類。
然則蒼等為術猶未足以博盡群數也。
徽尋九數有重差之名，原其指趣乃所以施於此也。
凡望極高、測絕深而兼知其遠者必用重差，句股則必以重差為率，故曰重差也。
立兩表於洛陽之城，令高八尺。南北各盡平地，同日度其正中之景。
以景差為法，表高乘表間為實，實如法而一，所得加表高，即日去地也。
以南表之景乘表間為實，實如法而一，即為從南表至南戴日下也。
以南戴日下及日去地為句、股，為之求弦，即日去人也。
以徑寸之筩南望日，日滿筩空，則定筩之長短以為股率，以筩徑為句率，
日去人之數為大股，大股之句即日徑也。
雖天圓穹之象猶曰可度，又況泰山之高與江海之廣哉。
徽以為今之史籍且略舉天地之物，考論厥數，載之於志，以闡世術之美。
輒造重差，并為注解，以究古人之意，綴於句股之下。
度高者重表，測深者累矩，孤離者三望，離而又旁求者四望。
觸類而長之，則雖幽遐詭伏，靡所不入。
\end{prefacetext}

\begin{flushright}
\textit{博物君子，詳而覽焉。}
\end{flushright}

\newpage


%% %%%%%%%%%%%%%%%%%%%%%%%%%%%%%%%%%%%%%%%%%%%%%%%%%%%%%%%%%%%%%%%%%%%%%%%%%%%%%
%% VOLUME 1: FANGTIAN
%% %%%%%%%%%%%%%%%%%%%%%%%%%%%%%%%%%%%%%%%%%%%%%%%%%%%%%%%%%%%%%%%%%%%%%%%%%%%%%
\section{卷第一：方田}
\label{sec:fangtian}
\addcontentsline{toc}{subsection}{Fangti\'{a}n -- Rectangular Fields}

\begin{center}
\textit{方田以御田疇界域}
\end{center}

\noindent\textbf{方田}者，田之正也。諸田不等，以方為正，故曰方田。
此章以御田疇界域，論各種田地面積之計算法，并及分數四則運算之法。
方田之術，乃算術之根本。凡田有廣從，廣者橫也，從者縱也。廣從相乘，得積步。
積步者，田之冪也。冪者，方田之面積也。以積步除之以畝法，即得畝數。
畝法二百四十步，此古法也。百畝為頃，此田疇之大數也。

\noindent 本章所論之田形，凡有十數：方田、里田、以廣從相乘為法；
又有圭田，半廣以乘正從；邪田、箕田，并兩廣而半之以乘正從；
圓田，半周半徑相乘；弧田，弦矢之法；環田，并中外周而半之以乘徑。
諸田之術，皆由方田而生，以盈虛相補，化曲為直，乃得積步。

\bigskip

\noindent\textbf{\large 一、方田術與里田術}
\vspace{0.5em}

\begin{problem}[一]
今有田廣十五步，從十六步。問為田幾何？
\end{problem}

\begin{answer}
一畝。
\end{answer}

\begin{problem}[二]
又有田廣十二步，從十四步。問為田幾何？
\end{problem}

\begin{answer}
一百六十八步。
\end{answer}

\begin{method}
方田術曰：廣從步數相乘得積步。以畝法二百四十步除之，即畝數。百畝為一頃。
\end{method}

\begin{liuhuicomm}
此積謂田之冪。凡廣從相乘謂之冪。畝法二百四十步者，古制也。
凡田冪之術，廣從相乘，此其正術。何則？假令有田，廣十步，從十步，相乘得一百步。
以畝法二百四十步約之，不足一畝。若廣二十步，從十二步，相乘得二百四十步，正合一畝。
此廣從相乘之理，猶方田之積，不可易也。
\end{liuhuicomm}

\begin{problem}[三]
今有田廣一里，從一里。問為田幾何？
\end{problem}

\begin{answer}
三頃七十五畝。
\end{answer}

\begin{problem}[四]
又有田廣二里，從三里。問為田幾何？
\end{problem}

\begin{answer}
二十二頃五十畝。
\end{answer}

\begin{method}
里田術曰：廣從里數相乘得積里。以三百七十五乘之，即畝數。
\end{method}

\begin{liuhuicomm}
按：一里田，廣從各一里。里法三百步，廣從相乘得九萬步。以畝法二百四十步除之，
得三百七十五畝，即三頃七十五畝。故里田術以三百七十五乘積里者，
蓋一里積九萬步，以二百四十除之，得三百七十五也。此術簡便，不必復以三百步自乘、
又以二百四十步除之也。
\end{liuhuicomm}

\vspace{0.8em}
\noindent\textbf{\large 二、約分術與分數運算}
\vspace{0.5em}

\noindent 方田之章，非獨論田冪而已，亦所以明分數四則之術也。
分數者，算之綱紀。物之數量不可悉全，必以分言之。
約分者，約之使簡；合分者，合之使一；減分者，減之使餘；
課分者，課之使知多寡；平分者，平之使均；經分者，經之使各得；
乘分者，乘之使相廣；大廣田者，帶分以相乘也。
凡此八術，分數之全體也。

\begin{problem}[五]
今有十八分之十二。問約之得幾何？
\end{problem}

\begin{answer}
三分之二。
\end{answer}

\begin{problem}[六]
又有九十一分之四十九。問約之得幾何？
\end{problem}

\begin{answer}
十三分之七。
\end{answer}

\begin{method}
約分術曰：可半者半之，不可半者，副置分母子之數，以少減多，更相減損，求其等也。以等數約之。
\end{method}

\begin{liuhuicomm}
約分者，物之數量不可悉全，必以分言之。等數者，約之篤也。
按：約分術者，求分子分母之等數，以等數約之。等數者何？最大公約數也。
更相減損者，輾轉相除之法也。以少減多，餘者更以減少，遞推之至餘數為零，
則最後之減數即等數也。如十八分之十二，副置十八、十二，以十二減十八，餘六；
又以六減十二，適盡。故六為等數。以六約十八得三，約十二得二，故為三分之二也。
又如九十一分之四十九，副置九十一、四十九，以四十九減九十一，餘四十二；
以四十二減四十九，餘七；以七減四十二，適盡。故七為等數。以七約九十一得十三，
約四十九得七，故為十三分之七也。
\end{liuhuicomm}

\begin{problem}[七]
今有三分之一，五分之二。問合之得幾何？
\end{problem}

\begin{answer}
十五分之十一。
\end{answer}

\begin{problem}[八]
又有三分之二，七分之四，九分之五。問合之得幾何？
\end{problem}

\begin{answer}
得一、六十三分之五十。
\end{answer}

\begin{problem}[九]
又有二分之一，三分之二，四分之三，五分之四。問合之得幾何？
\end{problem}

\begin{answer}
得二、六十分之四十三。
\end{answer}

\begin{method}
合分術曰：母互乘子，并為實，母相乘為法，實如法而一。不滿法者，以法命之。其母同者，直相從之。
\end{method}

\begin{liuhuicomm}
合分者，并數之義也。母互乘子，約而言之。母相乘為法者，齊而同之也。
按：合分術者，分數之加法也。母互乘子者，齊同之術也。分母不同，不可徑并，
故先齊而同之。齊者，以子互乘母，使諸分同母；同者，以母相乘為共母。
如三分之一與五分之二，三乘二得六，五乘一得五，并得十一為實；母三與五相乘得十五為法。
實如法而一，得十五分之十一。若母同者，如二分之一與二分之一，直相從之，得一也。
其三個以上分數，術亦如之。如三分之二、七分之四、九分之五，三母互乘各子，
七乘九乘二得一百二十六，三乘九乘四得一百零八，三乘七乘五得一百零五，
并得三百三十九為實；母三乘七乘九得一百八十九為法。實如法而一，得一、一百八十九分之一百五十，
約之得一、六十三分之五十也。
\end{liuhuicomm}

\begin{problem}[一０]
今有九分之八，減其五分之一。問餘幾何？
\end{problem}

\begin{answer}
四十五分之三十一。
\end{answer}

\begin{problem}[一一]
又有四分之三，減其三分之一。問餘幾何？
\end{problem}

\begin{answer}
十二分之五。
\end{answer}

\begin{method}
減分術曰：母互乘子，以少減多，餘為實，母相乘為法，實如法而一。
\end{method}

\begin{liuhuicomm}
減分者，分數之減法也。其術與合分相似，唯不并而相減為異耳。
母互乘子者，亦齊同之術也。以少減多，取其餘也。餘為實，母相乘為法，實如法而一，
即得減後之數。如九分之八減五分之一，母互乘子，九乘一得九，五乘八得四十，
以九減四十，餘三十一為實；母相乘得四十五為法。故得四十五分之三十一。
\end{liuhuicomm}

\begin{problem}[一二]
今有八分之五，二十五分之十六。問孰多？多幾何？
\end{problem}

\begin{answer}
二十五分之十六多，多二百分之三。
\end{answer}

\begin{problem}[一三]
又有九分之八，七分之六。問孰多？多幾何？
\end{problem}

\begin{answer}
九分之八多，多六十三分之二。
\end{answer}

\begin{problem}[一四]
又有二十一分之八，五十分之十七。問孰多？多幾何？
\end{problem}

\begin{answer}
二十一分之八多，多一千五十分之四十三。
\end{answer}

\begin{method}
課分術曰：母互乘子，以少減多，餘為實，母相乘為法，實如法而一，即相多也。
\end{method}

\begin{liuhuicomm}
課分者，課較兩分之多寡也。母互乘子，齊而同之，然後可以相減。餘者即相多之數也。
如八分之五與二十五分之十六，母互乘子，八乘十六得一百二十八，二十五乘五得一百二十五，
以一百二十五減一百二十八，餘三為實；母相乘得二百為法。故二十五分之十六多，
多二百分之三也。此術猶分數之減法，但以課多寡為目的，故別立課分之名。
\end{liuhuicomm}

\begin{problem}[一五]
今有三分之一，三分之二，四分之三。問減多益少，各幾何而平？
\end{problem}

\begin{answer}
減四分之三者二，三分之二者一，并以益三分之一，而各平於十二分之七。
\end{answer}

\begin{problem}[一六]
又有二分之一，三分之二，四分之三。問減多益少，各幾何而平？
\end{problem}

\begin{answer}
減三分之二者一，四分之三者四，并以益二分之一，而各平於三十六分之二十三。
\end{answer}

\begin{method}
平分術曰：母互乘子，副并為平實，母相乘為法。以列數乘未并者各自為列實。亦以列數乘法，以平實減列實，餘，約之為所減。并所減以益於少，以法命平實，各得其平。
\end{method}

\begin{liuhuicomm}
平分者，均分不平之數，使各相等也。母互乘子，副并為平實者，求其平均之數也。
以列數乘未并者各自為列實者，各以列數乘其未并之數，使可比較也。
平實減列實，餘者即多者所當減之數；并所減以益於少者，合多者之餘以補少者之不足也。
如三分之一、三分之二、四分之三，母互乘子得四、八、九，副并得二十一為平實。
母相乘得三十六為法。列數三乘未并者，三乘四得十二，三乘八得二十四，三乘九得二十七，
各自為列實。以列數三乘法三十六得一百零八，以平實二十一減列實：
二十四減二十一餘三，二十七減二十一餘六，約之得所減。并所減以益十二，各平於三十六分之七。
\end{liuhuicomm}

\begin{problem}[一七]
今有七人，分八錢三分錢之一。問人得幾何？
\end{problem}

\begin{answer}
人得一錢、二十一分錢之四。
\end{answer}

\begin{problem}[一八]
又有三人，三分人之一，分六錢三分錢之一，四分錢之三。問人得幾何？
\end{problem}

\begin{answer}
人得二錢、八分錢之一。
\end{answer}

\begin{method}
經分術曰：以人數為法，錢數為實，實如法而一。有分者通之，重有分者同而通之。
\end{method}

\begin{liuhuicomm}
經分者，徑以人數分錢之義。通分者，齊而同之，使可分也。
按：經分者，分數之除法也。以人數為法，錢數為實，實如法而一，即每人所得之數。
有分者通之者，若錢數或人數有分數，則先通分使為整數，然後可除也。
如七人分八錢三分錢之一，人數七為整數，錢數八錢三分錢之一即三分之二十五。
通之，以分母三乘八得二十四，加一得二十五，即三分之二十五。以人數七為法，
錢數二十五為實（人數七亦通分，乘三得二十一為法），實如法而一，得二十一分之二十五，
即一錢二十一分錢之四也。
\end{liuhuicomm}

\begin{problem}[一九]
今有田廣七分步之四，從五分步之三。問為田幾何？
\end{problem}

\begin{answer}
三十五分步之十二。
\end{answer}

\begin{problem}[二０]
又有田廣九分步之七，從十一分步之九。問為田幾何？
\end{problem}

\begin{answer}
十一分步之七。
\end{answer}

\begin{problem}[二一]
又有田廣五分步之四，從九分步之五，問為田幾何？
\end{problem}

\begin{answer}
九分步之四。
\end{answer}

\begin{method}
乘分術曰：母相乘為法，子相乘為實，實如法而一。
\end{method}

\begin{liuhuicomm}
乘分者，分數之乘法也。母相乘為法，子相乘為實，實如法而一，即得乘積之數。
按：乘分所以求分田之冪也。如廣七分步之四，從五分步之三，母相乘七乘五得三十五為法，
子相乘四乘三得十二為實。實如法而一，得三十五分步之十二，即田之積步也。
此術之理，猶整數廣從之相乘，但以分數為廣從，故以母子各相乘也。
\end{liuhuicomm}

\begin{problem}[二二]
今有田廣三步、三分步之一，從五步、五分步之二。問為田幾何？
\end{problem}

\begin{answer}
十八步。
\end{answer}

\begin{problem}[二三]
又有田廣七步、四分步之三，從十五步、九分步之五。問為田幾何？
\end{problem}

\begin{answer}
一百二十步、九分步之五。
\end{answer}

\begin{problem}[二四]
又有田廣十八步、七分步之五，從二十三步、十一分步之六。問為田幾何？
\end{problem}

\begin{answer}
一畝二百步、十一分步之七。
\end{answer}

\begin{method}
大廣田術曰：分母各乘其全，分子從之，相乘為實。分母相乘為法。實如法而一。
\end{method}

\begin{liuhuicomm}
大廣田者，帶分數之廣從相乘也。分母各乘其全，分子從之者，化帶分數為假分數也。
如廣三步、三分步之一，分母三乘全步三得九，加子一得十，即三分之十。
從五步、五分步之二，分母五乘全步五得二十五，加子二得二十七，即五分之二十七。
母子各相乘，十乘二十七得二百七十為實；分母相乘三乘五得十五為法。
實如法而一，得十八步。此術與乘分同，唯先化帶分為假分，然後相乘耳。
\end{liuhuicomm}

\vspace{0.8em}
\noindent\textbf{\large 三、圭田術}
\vspace{0.5em}

\noindent 圭田者，三角之田也。形如圭璧之半，故曰圭田。
其術半廣以乘正從，亦可以半正從以乘廣，其積一也。
何則？以盈補虛，化圭田為直田之半也。

\begin{problem}[二五]
今有圭田廣十二步，正從二十一步。問為田幾何？
\end{problem}

\begin{answer}
一百二十六步。
\end{answer}

\begin{problem}[二六]
又有圭田廣五步、二分步之一，從八步、三分步之二。問為田幾何？
\end{problem}

\begin{answer}
二十三步、六分步之五。
\end{answer}

\begin{method}
術曰：半廣以乘正從。
\end{method}

\begin{liuhuicomm}
半廣者，以盈補虛，為直田也。亦可以半正從以乘廣。
按：圭田者，三角形田也。其術半廣以乘正從者，以三角形之面積等于同底同高矩形之半也。
假令有圭田，廣十二步，正從二十一步。半廣得六步，以乘正從二十一步，
得一百二十六步，即圭田之積也。若以半正從乘廣，半正從得十步半，以乘廣十二步，
亦得一百二十六步，其理一也。何則？兩正三角可合成一矩形，故三角之積必為矩形之半。
以盈補虛之法，割三角之盈以補其虛，即可合成半矩形也。
\end{liuhuicomm}

\vspace{0.8em}
\noindent\textbf{\large 四、邪田術}
\vspace{0.5em}

\noindent 邪田者，兩邊不平之田也。一頭廣、一頭狹，形如梯形。
并兩邪而半之，猶求平均之廣，以乘正從，即得積步。

\begin{problem}[二七]
今有邪田，一頭廣三十步，一頭廣四十二步，正從六十四步。問為田幾何？
\end{problem}

\begin{answer}
九畝一百四十四步。
\end{answer}

\begin{problem}[二八]
又有邪田，正廣六十五步，一畔從一百步，一畔從七十二步。問為田幾何？
\end{problem}

\begin{answer}
二十三畝七十步。
\end{answer}

\begin{method}
術曰：并兩邪而半之，以乘正從若廣。又可半正從若廣，以乘并，畝法而一。
\end{method}

\begin{liuhuicomm}
邪田者，梯形之田也。并兩邪而半之者，猶求其平均之廣也。
按：邪田之形，兩邊平行而廣狹不等。并兩廣而半之，得平均廣，以乘正從，即得積步。
此術之理，猶以盈補虛，化邪田為直田也。割邪田之兩端，以其盈補其虛，即可合成矩形。
如一頭廣三十步，一頭廣四十二步，并之得七十二步，半之得三十六步。
以乘正從六十四步，得二千三百零四步。以畝法二百四十步除之，得九畝餘一百四十四步，
即九畝一百四十四步也。又可半正從，以乘并者，半正從得三十二步，乘七十二步，
亦得二千三百零四步，理同也。
\end{liuhuicomm}

\vspace{0.8em}
\noindent\textbf{\large 五、箕田術}
\vspace{0.5em}

\noindent 箕田者，形如簸箕之田也。舌廣而踵狹，亦梯形之類。
箕田之術，與邪田同。并踵舌而半之，以乘正從，畝法而一。

\begin{problem}[二九]
今有箕田，舌廣二十步，踵廣五步，正從三十步。問為田幾何？
\end{problem}

\begin{answer}
一畝一百三十五步。
\end{answer}

\begin{problem}[三０]
又有箕田，舌廣一百一十七步，踵廣五十步，正從一百三十五步。問為田幾何？
\end{problem}

\begin{answer}
四十六畝二百三十二步半。
\end{answer}

\begin{method}
術曰：并踵舌而半之，以乘正從。畝法而一。
\end{method}

\begin{liuhuicomm}
箕田者，形如箕。踵者，狹端也；舌者，廣端也。并而半之，猶兩邪之術也。
按：箕田即邪田之異名耳。舌為廣頭，踵為狹頭，與邪田之一頭廣、一頭狹無異。
故其術亦同：并踵舌而半之，以乘正從。如舌廣二十步，踵廣五步，并之得二十五步，
半之得十二步半。以乘正從三十步，得三百七十五步。以畝法二百四十步除之，
得一畝餘一百三十五步，即一畝一百三十五步也。
\end{liuhuicomm}

\vspace{0.8em}
\noindent\textbf{\large 六、圓田術}
\vspace{0.5em}

\noindent 圓田者，圓形之田也。圓者，徑一周三，此其大略也。
然徽以為圓周之率，非周三徑一所能盡。故立割圓之術，以六觚為始，
割之又割，以至於不可割，則與圓周合體而無所失矣。

\begin{problem}[三一]
今有圓田，周三十步，徑十步。問為田幾何？
\end{problem}

\begin{answer}
七十五步。
\end{answer}

\begin{problem}[三二]
又有圓田，周一百八十一步，徑六十步、三分步之一。問為田幾何？
\end{problem}

\begin{answer}
十一畝九十步、十二分步之一。
\end{answer}

\begin{method}
術曰：半周半徑相乘得積步。

又術曰：周徑相乘，四而一。

又術曰：徑自相乘，三之，四而一。

又術曰：周自相乘，十二而一。
\end{method}

\begin{liuhuicomm}
此于徽術，當為田十畝二百八步三百一十四分步之一百一。臣淳風等謹依密率，為田十畝二百步八十八分步之八十七。

按：半周為從，半徑為廣，故廣從相乘為積步也。假令圓徑二尺，圓中容六觚之一面，與圓徑之半，其數均等。合徑率一而外周率三也。

又按：為圖，以六觚之一面乘一弧半徑，三之，得十二觚之冪。若又割之，次以十二觚之一面乘一弧之半徑，六之，則得二十四觚之冪。割之彌細，所失彌少。割之又割，以至于不可割，則與圓周合體而無所失矣。觚面之外，又有餘徑。以面乘餘徑，則冪出觚表。若夫觚之細者，與圓合體，則表無餘徑。表無餘徑，則冪不外出矣。以一面乘半徑，觚而裁之，每輒自倍。故以半周乘半徑而為圓冪。

此一周、徑，謂至然之數，非周三徑一之率也。周三者，從其六觚之環耳。以推圓規多少之覺，乃弓之與弦也。然世傳此法，莫肯精核；學者踵古，習其謬失。不有明據，辯之斯難。凡物類形象，不圓則方。方圓之率，誠著于近，則雖遠可知也。由此言之，其用博矣。謹按圖驗，更造密率。

按：圓田之術，古法以周三徑一為率。半周乘半徑，周三十步半之得十五，徑十步半之得五，
相乘得七十五步，此古法也。然周三徑一，乃六觚之周率，非真圓之率也。
徽以割圓術推之，圓周之真率，當為徑一百一十三，周三百五十五，是為密率。
以密率計之，周一百八十一步者，徑當為六十一步、一百七十七分步之四十三，
非六十步三分步之一也。以古法徑六十步三分步之一計之，積十一畝九十步十二分步之一；
以徽術計之，當為十畝二百八步三百一十四分步之一百一，較之古法，所差甚巨。
故割圓之術，不可不察也。割圓術者，以圓內接正六邊形為始，謂之六觚。
六觚之周等于圓徑之三，此周三徑一所由來也。以次割之，十二觚、二十四觚、
四十八觚、九十六觚，割之彌多，觚之周彌近于圓周。以九十六觚之冪推之，
圓冪之率，得徑一百一十三，周三百五十五，此徽之密率也。
\end{liuhuicomm}

\begin{problem}[三三]
今有宛田，下周三十步，徑十六步。問為田幾何？
\end{problem}

\begin{answer}
一百二十步。
\end{answer}

\begin{problem}[三四]
又有宛田，下周九十九步，徑五十一步。問為田幾何？
\end{problem}

\begin{answer}
五畝六十二步、四分步之一。
\end{answer}

\begin{method}
術曰：以徑乘周，四而一。
\end{method}

\begin{liuhuicomm}
宛田者，中央隆高，形如覆碗。此術以徑乘周四而一，亦猶圓冪之義也。
按：宛田者，球冠形之田也。中央隆高，四畔低下，形如覆碗之狀。
其術以徑乘周，四而一者，猶圓冪之周徑相乘四而一也。
如下周三十步，徑十六步，相乘得四百八十步，四而一得一百二十步。
以畝法二百四十步除之，不足一畝，即一百二十步也。
然宛田之術，以徽術考之，未為精確。蓋宛田中央隆高，其表面非平面，
此術所算，乃其大畧耳。
\end{liuhuicomm}

\begin{problem}[三五]
今有弧田，弦三十步，矢十五步。問為田幾何？
\end{problem}

\begin{answer}
一畝九十七步半。
\end{answer}

\begin{problem}[三六]
又有弧田，弦七十八步、二分步之一，矢十三步、九分步之七。問為田幾何？
\end{problem}

\begin{answer}
二畝一百五十五步、八十一分步之五十六。
\end{answer}

\begin{method}
術曰：以弦乘矢，矢又自乘，并之，二而一。
\end{method}

\begin{liuhuicomm}
弧田者，圓之殘也。弦者，弧之對徑也；矢者，弧之中至弦之垂距也。
按：弧田者，弓形之田也。形如弓背，故謂之弧田。
弦者，弓弦也，即弧之兩端之直線距離。矢者，箭矢也，即弧之中點至弦之垂距。
其術以弦乘矢，矢又自乘，并之，二而一者，以弦矢之積為三角形之積，
矢自乘為正方形之積，并而半之，以為弧田之近似積也。
如弦三十步，矢十五步，弦乘矢得四百五十步，矢自乘得二百二十五步，
并得六百七十五步，二而一得三百三十七步半。
以畝法二百四十步除之，得一畝餘九十七步半，即一畝九十七步半也。
然此術所得，乃弧田之近似數。若矢大于弦之半，則所失漸多。
徽以為當以割圓之術，求弧之周與矢之關係，乃得精確之數。
\end{liuhuicomm}

\begin{problem}[三七]
今有環田，中周九十二步，外周一百二十二步，徑五步。問為田幾何？
\end{problem}

\begin{answer}
二畝五十五步。
\end{answer}

\begin{problem}[三八]
又有環田，中周六十二步、四分步之三，外周一百一十三步、二分步之一，徑十二步、三分步之二。問為田幾何？
\end{problem}

\begin{answer}
四畝一百五十六步、四分步之一。
\end{answer}

\begin{method}
術曰：并中外周而半之，以徑乘之為積步。

密率術曰：置中外周步數，分母、子各居其下。母互乘子，通全步，內分子。以中周減外周，餘半之，以益中周。徑亦通分內子，以乘周為實。分母相乘為法，除之為積步，餘積步之分。以畝法除之，即畝數也。
\end{method}

\begin{liuhuicomm}
環田者，如圓環之形。并中外周而半之者，猶并兩邪而半之也。
按：環田者，圓環形之田也。外周為大圓之周，中周為小圓之周，徑為兩圓周之距離。
并中外周而半之者，猶求平均之周也。以徑乘平均周，即得環田之積。
如中周九十二步，外周一百二十二步，并之得二百一十四步，半之得一百零七步。
以徑五步乘之，得五百三十五步。以畝法二百四十步除之，得二畝餘五十五步，
即二畝五十五步也。此術之理，猶以環田展開為直條，平均周為長，徑為廣，
廣從相乘得積步也。若以密率術算之，先通分，母互乘子，內分子，
以中周減外周，餘半之，以益中周，得平均周。徑亦通分，以乘平均周為實，
分母相乘為法，實如法而一，即積步也。密率術所以處理帶分數之環田也。
\end{liuhuicomm}

\vspace{0.8em}
\noindent\textbf{\large 七、方田章小結}
\vspace{0.5em}

\begin{quoting}
方田之章，凡三十八問。論田冪之計算，凡有八術：方田、里田、圭田、邪田、
箕田、圓田、弧田、環田。又論分數之運算，凡有八術：約分、合分、減分、
課分、平分、經分、乘分、大廣田。此十六章，皆算術之基礎也。

劉徽之注，尤為精要。其論圓田也，創割圓之術，以六觚為始，割之又割，
以至不可割，則與圓周合體。此極限思想之濫觴也。其論分數也，以齊同為綱，
母互乘子，母相乘為法，條理井然。方田之術，上承商高，下啟祖沖之，
為中國古代數學之瑰寶也。
\end{quoting}

\newpage


%% %%%%%%%%%%%%%%%%%%%%%%%%%%%%%%%%%%%%%%%%%%%%%%%%%%%%%%%%%%%%%%%%%%%%%%%%%%%%%
%% VOLUME 2: SUMI
%% %%%%%%%%%%%%%%%%%%%%%%%%%%%%%%%%%%%%%%%%%%%%%%%%%%%%%%%%%%%%%%%%%%%%%%%%%%%%%
\section{卷第二：粟米}
\label{sec:sumi}
\addcontentsline{toc}{subsection}{S\`{u}m\u{i} -- Millet and Rice}

\begin{center}
\textit{粟米以御交質變易}
\end{center}

\noindent 本章論各種谷物之交換比率，及比例計算之法。
\textit{今有術}（即今之三率法、比例法）為本章之核心算法。

\noindent 粟米者，五穀之總名也。粟為百穀之長，故以為率之標準。
本章先列粟米之法，凡二十種谷物，各以其率，然後以今有術求之。
今有術者，以所有數乘所求率為實，以所有率為法，實如法而一。
此術簡便，而用之無窮，比例計算之總綱也。

\bigskip

\begin{center}
\textbf{\large 粟米之法}
\end{center}

\vspace{0.5em}

\noindent
\begin{tabular}{llll}
\textbf{粟率五十} & \textbf{糲米三十} & \textbf{粺米二十七} & \textbf{鑿米二十四} \\
\textbf{御米二十一} & 小\textbullet{}十三半 & 大\textbullet{}五十四 & \textbf{糲飯七十五} \\
\textbf{粺飯五十四} & \textbf{鑿飯四十八} & \textbf{御飯四十二} & \textbf{菽、荅、麻、麥各四十五} \\
\textbf{稻六十} & \textbf{豉六十三} & \textbf{飧九十} & \textbf{熟菽一百三半} \\
\multicolumn{4}{l}{\textbf{櫱一百七十五}} \\
\end{tabular}

\vspace{0.5em}

\begin{method}
今有術曰：以所有數乘所求率為實，以所有率為法，實如法而一。
\end{method}

\begin{liuhuicomm}
此術今有之義，凡數相與者謂之率。率者，自相與通。有分則散，無分則不散。凡此諸率，各有所求，故曰今有術。

按：今有術者，比例之術也。以所有數乘所求率為實，以所有率為法，實如法而一，
即得所求之數。何謂率？率者，兩數相比之關係也。如粟率五十、糲米三十，
即五十粟可換三十糲米也。今有術之義，猶言：「今有粟若干，欲為糲米，問得幾何？」
故以所有粟數乘所求糲米率三十為實，以所有粟率五十為法，實如法而一，
即得糲米之數。此術之所以謂之「今有」者，蓋「今有」即「現有」之義，
謂現有之物數，欲換為他物，故以比例求之也。

率者，又為分數之綱紀。凡數相與者謂之率，率知，則可以互求。
有分則散者，若率中有分數，則散而通之，使為整數。無分則不散者，
率中無分數，則徑以整數相求，不必更散也。此今有術之大要也。
\end{liuhuicomm}

\vspace{0.8em}
\noindent\textbf{\large 一、以粟求米}
\vspace{0.5em}

\noindent 粟為五穀之長，故以粟為基準。凡以粟求各種米飯，皆用今有術。
粟率五十，為所有率；所求米之率，為所求率；所有粟數，為所有數。

\begin{problem}[一]
今有粟一斗，欲為糲米。問得幾何？
\end{problem}

\begin{answer}
為糲米六升。
\end{answer}

\begin{method}
術曰：以粟求糲米，三之，五而一。
\end{method}

\begin{liuhuicomm}
按：粟率五十，糲米率三十。以粟一斗乘糲米率三十，得三十，為實。
以粟率五十為法。實如法而一，得五分之三斗，即六升也。
術云「三之，五而一」者，即三十乘之、五十除之之簡語也。約而言之，
三十之於五十，即三之於五，故云三之、五而一也。此約率之術，使計算簡便也。
\end{liuhuicomm}

\begin{problem}[二]
今有粟二斗一升，欲為粺米。問得幾何？
\end{problem}

\begin{answer}
為粺米一斗一升、五十分升之十七。
\end{answer}

\begin{method}
術曰：以粟求粺米，二十七之，五十而一。
\end{method}

\begin{liuhuicomm}
按：粟率五十，粺米率二十七。以粟二斗一升乘粺米率二十七。
二斗一升即二十一升。二十一乘二十七得五百六十七，為實。
以粟率五十為法。實如法而一，得一十一升餘十七，即一斗一升、五十分升之十七也。
術云「二十七之，五十而一」者，即約率也。
\end{liuhuicomm}

\begin{problem}[三]
今有粟四斗五升，欲為鑿米。問得幾何？
\end{problem}

\begin{answer}
為鑿米二斗一升、五分升之三。
\end{answer}

\begin{method}
術曰：以粟求鑿米，十二之，二十五而一。
\end{method}

\begin{liuhuicomm}
按：粟率五十，鑿米率二十四。十二之、二十五而一者，約率也。
二十四之於五十，即十二之於二十五。以粟四斗五升，即四十五升，乘十二得五百四十，
為實。以二十五為法。實如法而一，得二十一升餘十五，即二斗一升五分升之三。
\end{liuhuicomm}

\begin{problem}[四]
今有粟七斗九升，欲為御米。問得幾何？
\end{problem}

\begin{answer}
為御米三斗三升、五十分升之九。
\end{answer}

\begin{method}
術曰：以粟求御米，二十一之，五十而一。
\end{method}

\begin{problem}[五]
今有粟一斗，欲為小●。問得幾何？
\end{problem}

\begin{answer}
為小●二升、一十分升之七。
\end{answer}

\begin{method}
術曰：以粟求小●，二十七之，百而一。
\end{method}

\begin{liuhuicomm}
按：小\textbullet{}率十三半，即二分之二十七也。
以粟一斗即十升，乘小\textbullet{}率二十七，得二百七十，為實。
以粟率五十乘二（因小\textbullet{}率有半，故通分）得一百，為法。
實如法而一，得二升餘七十，約之得二升一十分升之七。
術云「二十七之，百而一」者，即此之約率也。
\end{liuhuicomm}

\begin{problem}[六]
今有粟九斗八升，欲為大●。問得幾何？
\end{problem}

\begin{answer}
為大●一十斗五升、二十五分升之二十一。
\end{answer}

\begin{method}
術曰：以粟求大●，二十七之，二十五而一。
\end{method}

\begin{problem}[七]
今有粟二斗三升，欲為糲飯。問得幾何？
\end{problem}

\begin{answer}
為糲飯三斗四升半。
\end{answer}

\begin{method}
術曰：以粟求糲飯，三之，二而一。
\end{method}

\begin{liuhuicomm}
按：粟率五十，糲飯率七十五。三之二而一者，七十五之於五十，即三之於二也。
以粟二十三升乘三得六十九，為實。以二為法。實如法而一，得三十四升半，
即三斗四升半也。
\end{liuhuicomm}

\begin{problem}[八]
今有粟三斗六升，欲為粺飯。問得幾何？
\end{problem}

\begin{answer}
為粺飯三斗八升、二十五分升之二十二。
\end{answer}

\begin{method}
術曰：以粟求粺飯，二十七之，二十五而一。
\end{method}

\begin{problem}[九]
今有粟八斗六升，欲為鑿飯。問得幾何？
\end{problem}

\begin{answer}
為鑿飯八斗二升、二十五分升之一十四。
\end{answer}

\begin{method}
術曰：以粟求鑿飯，二十四之，二十五而一。
\end{method}

\begin{problem}[一０]
今有粟九斗八升，欲為御飯。問得幾何？
\end{problem}

\begin{answer}
為御飯八斗二升、二十五分升之八。
\end{answer}

\begin{method}
術曰：以粟求御飯，二十一之，二十五而一。
\end{method}

\begin{problem}[一一]
今有粟三斗少半升，欲為菽。問得幾何？
\end{problem}

\begin{answer}
為菽二斗七升、一十分升之三。
\end{answer}

\begin{problem}[一二]
今有粟四斗一升、太半升，欲為荅。問得幾何？
\end{problem}

\begin{answer}
為荅三斗七升半。
\end{answer}

\begin{problem}[一三]
今有粟五斗、太半升，欲為麻。問得幾何？
\end{problem}

\begin{answer}
為麻四斗五升、五分升之三。
\end{answer}

\begin{problem}[一四]
今有粟一十斗八升、五分升之二，欲為麥。問得幾何？
\end{problem}

\begin{answer}
為麥九斗七升、二十五分升之一十四。
\end{answer}

\begin{method}
術曰：以粟求菽、荅、麻、麥，皆九之，十而一。
\end{method}

\begin{liuhuicomm}
菽、荅、麻、麥之率各四十五，以九之十而一者，約而言之也。
按：菽、荅、麻、麥四物，率各四十五。四十五之於五十，即九之於十。
故以粟求此四物，皆九之十而一也。以粟數乘九，以十除之，即得各物之數。
\end{liuhuicomm}

\begin{problem}[一五]
今有粟七斗五升、七分升之四，欲為稻。問得幾何？
\end{problem}

\begin{answer}
為稻九斗、三十五分升之二十四。
\end{answer}

\begin{method}
術曰：以粟求稻，六之，五而一。
\end{method}

\begin{liuhuicomm}
按：粟率五十，稻率六十。六之五而一者，六十之於五十，即六之於五也。
以粟求稻，稻之率大于粟率，故所得多于原粟。以粟七斗五升七分升之四，
通分為七分之五百二十九，乘六得三千一百七十四，為實。以五乘七得三十五為法。
實如法而一，得九斗三十五分升之二十四也。
\end{liuhuicomm}

\begin{problem}[一六]
今有粟七斗八升，欲為豉。問得幾何？
\end{problem}

\begin{answer}
為豉九斗八升、二十五分升之七。
\end{answer}

\begin{method}
術曰：以粟求豉，六十三之，五十而一。
\end{method}

\begin{problem}[一七]
今有粟五斗五升，欲為飧。問得幾何？
\end{problem}

\begin{answer}
為飧九斗九升。
\end{answer}

\begin{method}
術曰：以粟求飧，九之，五而一。
\end{method}

\begin{liuhuicomm}
按：粟率五十，飧率九十。九之五而一者，九十之於五十，即九之於五也。
以粟五斗五升即五十五升，乘九得四百九十五，為實。以五為法。
實如法而一，得九十九升，即九斗九升也。
\end{liuhuicomm}

\begin{problem}[一八]
今有粟四斗，欲為熟菽。問得幾何？
\end{problem}

\begin{answer}
為熟菽八斗二升、五分升之四。
\end{answer}

\begin{method}
術曰：以粟求熟菽，二百七之，百而一。
\end{method}

\begin{liuhuicomm}
按：熟菽率一百三半，即二分之二百七也。
以粟四斗即四十升，乘二百七得八千二百八十，為實。
以粟率五十乘二（通分）得一百，為法。
實如法而一，得八十二升五分升之四，即八斗二升五分升之四也。
\end{liuhuicomm}

\begin{problem}[一九]
今有粟二斗，欲為櫱。問得幾何？
\end{problem}

\begin{answer}
為櫱七斗。
\end{answer}

\begin{method}
術曰：以粟求櫱，七之，二而一。
\end{method}

\begin{liuhuicomm}
按：粟率五十，櫱率一百七十五。七之二而一者，一百七十五之於五十，
約之即七之於二也。以粟二斗即二十升，乘七得一百四十，為實。
以二為法。實如法而一，得七十升，即七斗也。
\end{liuhuicomm}

\vspace{0.8em}
\noindent\textbf{\large 二、以米求粟（返求）}
\vspace{0.5em}

\noindent 既有以粟求米，亦有以米返求於粟。其術亦然，唯所求率與所有率相易耳。
以米求粟，則米之率為所有率，粟之率為所求率也。

\begin{problem}[二０]
今有糲米十五斗五升、五分升之二，欲為粟。問得幾何？
\end{problem}

\begin{answer}
為粟二十五斗九升。
\end{answer}

\begin{method}
術曰：以糲米求粟，五之，三而一。
\end{method}

\begin{liuhuicomm}
按：以糲米求粟，糲米率三十為所有率，粟率五十為所求率。
五十之於三十，即五之於三。以糲米十五斗五升五分升之二，
通分為五分之七十八，乘五得三百九十，為實。以三乘五得十五為法。
實如法而一，得二十六升，即二十五斗九升也。此返求之法，
與以粟求糲米相反，而其術則同。
\end{liuhuicomm}

\begin{problem}[二一]
今有粺米二斗，欲為粟。問得幾何？
\end{problem}

\begin{answer}
為粟三斗七升、二十七分升之一。
\end{answer}

\begin{method}
術曰：以粺米求粟，五十之，二十七而一。
\end{method}

\begin{problem}[二二]
今有鑿米三斗、少半升，欲為粟。問得幾何？
\end{problem}

\begin{answer}
為粟六斗三升、三十六分升之七。
\end{answer}

\begin{method}
術曰：以鑿米求粟，二十五之，十三而一。
\end{method}

\begin{problem}[二三]
今有御米十四斗，欲為粟。問得幾何？
\end{problem}

\begin{answer}
為粟三十三斗三升、少半升。
\end{answer}

\begin{method}
術曰：以御米求粟，五十之，二十一而一。
\end{method}

\begin{problem}[二四]
今有稻一十二斗六升、一十五分升之一十四，欲為粟。問得幾何？
\end{problem}

\begin{answer}
為粟一十斗五升、九分升之七。
\end{answer}

\begin{method}
術曰：以稻求粟，五之，六而一。
\end{method}

\begin{liuhuicomm}
按：以稻求粟，稻率六十為所有率，粟率五十為所求率。
五十之於六十，即五之於六。以稻十二斗六升十五分升之十四，
通分為十五分之一百九十四，乘五得九百七十，為實。
以六乘十五得九十為法。實如法而一，得一十斗五升九分升之七也。
\end{liuhuicomm}

\vspace{0.8em}
\noindent\textbf{\large 三、米飯互求}
\vspace{0.5em}

\noindent 粟米之章，非獨以粟求米、以米求粟而已，亦論各種米飯之間互相求換之法。
其術皆用今有術，以所求物之率為所求率，以所有物之率為所有率也。

\begin{problem}[二五]
今有糲米一十九斗二升、七分升之一，欲為粺米。問得幾何？
\end{problem}

\begin{answer}
為粺米一十七斗二升、一十四分升之一十三。
\end{answer}

\begin{method}
術曰：以糲米求粺米，九之，十而一。
\end{method}

\begin{liuhuicomm}
按：糲米率三十，粺米率二十七。二十七之於三十，即九之於十。
以糲米求粺米，所得當少于原糲米，以其率小也。
以糲米十九斗二升七分升之一，通分為七分之一千三百四十五，
乘九得一萬二千一百零五，為實。以十乘七得七十為法。
實如法而一，得一十七斗二升一十四分升之一十三也。
\end{liuhuicomm}

\begin{problem}[二六]
今有糲米六斗四升、五分升之三，欲為糲飯。問得幾何？
\end{problem}

\begin{answer}
為糲飯一十六斗一升半。
\end{answer}

\begin{method}
術曰：以糲米求糲飯，五之，二而一。
\end{method}

\begin{liuhuicomm}
按：糲米率三十，糲飯率七十五。七十五之於三十，即五之於二。
以糲米求糲飯，所得當多于原糲米一倍半，以飯率大于米率也。
蓋飯者，米之蒸熟者也。米蒸熟則膨脹，故其率大于米。
以糲米六斗四升五分升之三，通分為五分之三百二十三，
乘五得一千六百一十五，為實。以二乘五得十為法。
實如法而一，得一十六斗一升半也。
\end{liuhuicomm}

\begin{problem}[二七]
今有糲飯七斗六升、七分升之四，欲為飧。問得幾何？
\end{problem}

\begin{answer}
為飧九斗一升、三十五分升之三十一。
\end{answer}

\begin{method}
術曰：以糲飯求飧，六之，五而一。
\end{method}

\begin{problem}[二八]
今有菽一斗，欲為熟菽。問得幾何？
\end{problem}

\begin{answer}
為熟菽二斗三升。
\end{answer}

\begin{method}
術曰：以菽求熟菽，二十三之，十而一。
\end{method}

\begin{liuhuicomm}
按：菽率四十五，熟菽率一百三半，即二分之二百七也。
二百七之於四十五，即二十三之於十（約去公因數九）。
以菽一斗即十升，乘二十三得二百三十，為實。以十為法。
實如法而一，得二十三升，即二斗三升也。
\end{liuhuicomm}

\begin{problem}[二九]
今有菽二斗，欲為豉。問得幾何？
\end{problem}

\begin{answer}
為豉二斗八升。
\end{answer}

\begin{method}
術曰：以菽求豉，七之，五而一。
\end{method}

\begin{liuhuicomm}
按：菽率四十五，豉率六十三。六十三之於四十五，即七之於五（約去公因數九）。
以菽二斗即二十升，乘七得一百四十，為實。以五為法。
實如法而一，得二十八升，即二斗八升也。
\end{liuhuicomm}

\begin{problem}[三０]
今有麥八斗六升、七分升之三，欲為小●，問得幾何？
\end{problem}

\begin{answer}
為小●二斗五升、一十四分升之一十三。
\end{answer}

\begin{method}
術曰：以麥求小●，三之，十而一。
\end{method}

\begin{problem}[三一]
今有麥一斗，欲為大●。問得幾何？
\end{problem}

\begin{answer}
為大●一斗二升。
\end{answer}

\begin{method}
術曰：以麥求大●，六之，五而一。
\end{method}

\vspace{0.8em}
\noindent\textbf{\large 四、經率術與經術}
\vspace{0.5em}

\noindent 經率術者，以錢買物，求物之單價也。以所買率為法，所出錢數為實，
實如法得一錢。經術者，以錢買物，求物之總數也。以所求率乘錢數為實，
以所買率為法，實如法得一。二術相為表里，皆今有術之變化也。

\begin{problem}[三二]
今有出錢一百六十，買瓴甓十八枚。問枚幾何？
\end{problem}

\begin{answer}
一枚，八錢、九分錢之八。
\end{answer}

\begin{method}
經率術曰：以所買率為法，所出錢數為實，實如法得一錢。
\end{method}

\begin{liuhuicomm}
按：經率術者，求物之單價也。以所買率為法，所出錢數為實，
實如法而一，即每物之值錢數。以出錢一百六十，買瓴甓十八枚，
以十八為法，一百六十為實。實如法而一，得八錢餘十六。
以法十八命之，即八錢九分錢之八也。
\end{liuhuicomm}

\begin{problem}[三三]
今有出錢一萬三千五百，買竹二千三百五十箇。問箇幾何？
\end{problem}

\begin{answer}
一箇，五錢、四十七分錢之三十五。
\end{answer}

\begin{problem}[三四]
今有出錢五千七百八十五，買漆一斛六斗七升、太半升。欲斗率之，問斗幾何？
\end{problem}

\begin{answer}
一斗，三百四十五錢、五百三分錢之一十五。
\end{answer}

\begin{problem}[三五]
今有出錢七百二十，買縑一匹二丈一尺。欲丈率之，問丈幾何？
\end{problem}

\begin{answer}
一丈，一百一十八錢、六十一分錢之二。
\end{answer}

\begin{problem}[三六]
今有出錢二千三百七十，買布九匹二丈七尺。欲匹率之，問匹幾何？
\end{problem}

\begin{answer}
一匹，二百四十四錢、一百二十九分錢之一百二十四。
\end{answer}

\begin{problem}[三七]
今有出錢一萬三千六百七十，買絲一石二鈞一十七斤。欲石率之，問石幾何？
\end{problem}

\begin{answer}
一石，八千三百二十六錢、一百九十七分錢之一百七十八。
\end{answer}

\begin{method}
經術術曰：以所求率乘錢數為實，以所買率為法，實如法得一。
\end{method}

\begin{liuhuicomm}
按：經術者，求每單位物之值也。以所求率乘錢數為實，以所買率為法，
實如法而一，即每單位物之值。如出錢五千七百八十五，買漆一斛六斗七升太半升。
一斛六斗七升太半升，即一斛又六斗七升又二分之一升也。
以所求率一斗乘錢數五千七百八十五，得五千七百八十五，為實。
以所買率十六斗七升半，即十六斗又十五分斗之五，亦即二分之三十三斗，
通分為一百六十七升半，即二分之三百三十五升為法。
實如法而一，得三百四十五錢五百三分錢之十五也。
\end{liuhuicomm}

\vspace{0.8em}
\noindent\textbf{\large 五、其率術與反其率術}
\vspace{0.5em}

\noindent 其率術者，貴賤相雜之率也。以錢買物，物有貴賤二品，問各得幾何。
其術以所買之數為法，以錢數乘所求率為實。不滿法者反以實減法，法賤實貴。
反其率術者，與其率術相反也。以錢數為法，所率為實。

\begin{problem}[三八]
今有出錢五百七十六，買竹七十八箇。欲其大小率之，問各幾何？
\end{problem}

\begin{answer}
其四十八箇，箇七錢。其三十箇，箇八錢。
\end{answer}

\begin{method}
其率術曰：各置所買石、鈞、斤、兩以為法，以所率乘錢數為實，實如法而一。不滿法者反以實減法，法賤實貴。
\end{method}

\begin{liuhuicomm}
其率者，所求之率也。以所買之數為法，以錢數乘所求率為實，實如法而一。
不滿法者，反以實減法，法賤實貴，謂法為賤率，實為貴率也。
按：其率術者，貴賤相雜之算法也。如出錢五百七十六，買竹七十八箇，
箇有大小二品。以大箇箇八錢、小箇箇七錢為率，列之。
以所買七十八箇為法，以錢數乘所求率。若先求小箇之數，
以小箇率七乘錢數五百七十六得四千零三十二為實。
實如法而一，以七十八為法，得五十一餘五十四。不滿法，
反以實減法，七十八減五十四得二十四，即小箇數也。
反以五十四減七十八得二十四，為大箇數也。
法賤實貴者，七為小箇之率，即賤率；餘數為貴率之減數也。
\end{liuhuicomm}

\begin{problem}[三九]
今有出錢一千一百二十，買絲一石二鈞十八斤。欲其貴賤斤率之，問各幾何？
\end{problem}

\begin{answer}
其二鈞八斤，斤五錢。其一石一十斤，斤六錢。
\end{answer}

\begin{problem}[四０]
今有出錢一萬三千九百七十，買絲一石二鈞二十八斤三兩五銖。欲其貴賤石率之，問各幾何？
\end{problem}

\begin{answer}
其一石，八千三百二十六錢、一百九十七分錢之一百七十八。其二鈞二十八斤三兩五銖，五千六百三十九錢、一百九十七分錢之六十七。
\end{answer}

\begin{problem}[四一]
今有出錢一萬三千九百七十，買絲一石二鈞二十八斤三兩五銖。欲其貴賤鈞率之，問各幾何？
\end{problem}

\begin{answer}
其七鈞，斤二百六十一錢、一百九十七分錢之一百三十三。其一鈞，斤二百七十八錢、一百九十七分錢之一百三十二。
\end{answer}

\begin{problem}[四二]
今有出錢一萬三千九百七十，買絲一石二鈞二十八斤三兩五銖。欲其貴賤斤率之，問各幾何？
\end{problem}

\begin{answer}
其一石二鈞八斤，斤五錢。其一石一十斤，斤六錢。
\end{answer}

\begin{problem}[四三]
今有出錢一萬三千九百七十，買絲一石二鈞二十八斤三兩五銖。欲其貴賤兩率之，問各幾何？
\end{problem}

\begin{answer}
其一石一鈞一十七斤一十四兩，兩四錢。其一鈞一十斤五兩，兩五錢。
\end{answer}

\begin{problem}[四四]
今有出錢一萬三千九百七十，買絲一石二鈞二十八斤三兩五銖。欲其貴賤銖率之，問各幾何？
\end{problem}

\begin{answer}
其一石一鈞一十七斤一十四兩一銖，銖四錢。其一鈞一十斤五兩四銖，銖五錢。
\end{answer}

\begin{problem}[四五]
今有出錢六百一十，買羽二千一百翭。欲其貴賤率之，問各幾何？
\end{problem}

\begin{answer}
其一千一百四十翭，三翭一錢。其九百六十翭，四翭一錢。
\end{answer}

\begin{problem}[四六]
今有出錢九百八十，買矢簳五千八百二十枚。欲其貴賤率之，問各幾何？
\end{problem}

\begin{answer}
其三百枚，五枚一錢。其五千五百二十枚，六枚一錢。
\end{answer}

\begin{method}
反其率術曰：以錢數為法，所率為實，實如法而一。不滿法者反以實減法，法少，實多。二物各以所得多少之數乘法實，即物數。
\end{method}

\begin{liuhuicomm}
反其率者，與其率相反也。其率以所買為法，此以錢數為法，故曰反其率。
不滿法者反以實減法，法少實多，二物各以多少之數乘法實，即得各物之數也。
按：反其率術者，與其率術相反而相成也。其率術以所買物數為法，
反其率術以所出錢數為法。其率術求物之貴賤單價，反其率術求物之貴賤數量。
以錢數為法，所率為實，實如法而一。不滿法者，反以實減法，
法少實多，二物各以所得多少之數乘法實，即得各物之數。
如出錢六百一十，買羽二千一百翭。以錢數六百一十為法，
所買羽數二千一百為實。實如法而一，得三餘二百七十。
不滿法，反以實減法，六百一十減二百七十得三百四十，
即三翭一錢之羽數也。又以二百七十為四翭一錢之羽數也。
二物各以多少之數乘法實：三乘三百四十得一千零二十，
以錢除之得三百四十翭（此為三翭一錢之總數）；
四乘二百七十得一千零八十，以錢除之得二百七十翭（此為四翭一錢之總數）。
合之得二千一百翭也。
\end{liuhuicomm}

\vspace{0.8em}
\noindent\textbf{\large 六、粟米章小結}
\vspace{0.5em}

\begin{quoting}
粟米之章，凡四十六問。先列粟米之法二十種，各定其率。
然後以今有術求之。今有術者，以所有數乘所求率為實，以所有率為法，
實如法而一。此術為比例計算之總綱，不獨粟米為然，凡數之相與者皆可用之。

本章又有經率術、經術、其率術、反其率術，皆今有術之變化也。
經率術求單價，經術求單位值，其率術求貴賤之率，反其率術求貴賤之數。
四術相為表里，以御交質變易之繁。

劉徽之注，闡明率之義，謂「凡數相與者謂之率，率者自相與通」，
此比例理論之精髓也。李淳風等又按密率校之，使精度益進。
粟米之術，於今之商業計算、貨幣交換、物價換算，無不可用，其義博矣。
\end{quoting}

\newpage


%% %%%%%%%%%%%%%%%%%%%%%%%%%%%%%%%%%%%%%%%%%%%%%%%%%%%%%%%%%%%%%%%%%%%%%%%%%%%%%
%% VOLUME 3: CUIFEN
%% %%%%%%%%%%%%%%%%%%%%%%%%%%%%%%%%%%%%%%%%%%%%%%%%%%%%%%%%%%%%%%%%%%%%%%%%%%%%%
\section{卷第三：衰分}
\label{sec:cuifen}
\addcontentsline{toc}{subsection}{Cu\={i}f\={e}n -- Distributing by Proportion}

\begin{center}
\textit{衰分以御貴賤粟稅之率}
\end{center}

\noindent 本章論比例分配之法，包括按等級、按數量、按比率之分配問題。
\textit{衰分術}為本章之核心算法。

\noindent 衰分者，差分也。衰者，等差也。以各人之等差為率，分配財物或賦役，
使各得其分。衰分術曰：各置列衰，副并為法，以所分乘未并者各自為實，實如法而一。
不滿法者，以法命之。此術與今有術相通，唯以列衰為率，是其特點也。

\noindent 列衰者，列其差數以為率也。如爵位之高下、資產之多少、出力之大小，
皆可定為列衰。列衰既定，以副并為法，各以所分乘未并者為實，實如法而一，
即各得其分也。

\noindent 本章又有返衰術。返衰者，反其列衰也。高爵出少，以次漸多，
與正衰相反。返衰術曰：列置衰而令相乘，動者為不動者衰。

\bigskip

\begin{method}
衰分術曰：各置列衰，副并為法，以所分乘未并者各自為實，實如法而一。不滿法者，以法命之。
\end{method}

\begin{liuhuicomm}
衰分者，差分也。列衰者，列其差數以為率也。各置列衰，副并為法，以所分乘未并者各自為實，實如法而一，即各得其分也。

按：衰分之術，比例分配之法也。衰者，等級之差也。如大夫五、不更四、簪裹三、上造二、公士一，此爵位之衰也。各置列衰者，列各人之等級差數。副并為法者，以各衰數相加為法。以所分之物乘未并之各衰數為實，實如法而一，即各人所分得之物數也。

其理與今有術相通。今有術以所有率為法，以所求率乘所有數為實。衰分術以列衰之并為法，以各列衰乘所分數為實。所不同者，今有術只有一個所求率，衰分術有多個列衰耳。

不滿法者以法命之者，若實小于法，則以分數命之，使各得其零數也。此與經分之意同。
\end{liuhuicomm}

\vspace{0.8em}
\noindent\textbf{\large 一、以爵次分物}
\vspace{0.5em}

\noindent 以爵次分物者，按爵位之高下列衰分配也。秦漢之制，爵有二十級，
本章所列有大夫、不更、簪裹、上造、公士五等。大夫第五級，不更第四級，
簪裹第三級，上造第二級，公士第一級。以爵級為列衰，級高者分多，級低者分少。

\begin{problem}[一]
今有大夫、不更、簪裹、上造、公士，凡五人，共獵得五鹿。欲以爵次分之，問各得幾何？
\end{problem}

\begin{answer}
大夫得一鹿、三分鹿之二。不更得一鹿、三分鹿之一。簪裹得一鹿。上造得三分鹿之二。公士得三分鹿之一。
\end{answer}

\begin{method}
術曰：列置爵數，各自為衰，副并為法。以五鹿乘未并者，各自為實。實如法得一鹿。
\end{method}

\begin{liuhuicomm}
大夫五，不更四，簪裹三，上造二，公士一，是為列衰。副并得一十五為法。以五鹿乘未并者，大夫得二十五，不更得二十，簪裹得一十五，上造得一十，公士得五，各自為實。實如法而一，即各得鹿數也。

按：此術列置爵數為衰者，以爵級為列衰也。大夫五級，不更四級，簪裹三級，上造二級，公士一級，故以五、四、三、二、一為列衰。副并得十五為法。以所分五鹿乘未并之各列衰，大夫得五乘五二十五，不更得五乘四二十，簪裹得五乘三十五，上造得五乘二十，公士得五乘一五，各自為實。以法十五除各實：大夫二十五除以十五得一又三分之二，即一鹿三分鹿之二；不更二十除以十五得一又三分之一，即一鹿三分鹿之一；簪裹十五除以十五得一鹿；上造十除以十五得三分之二鹿；公士五除以十五得三分之一鹿。此衰分之正法也。
\end{liuhuicomm}

\vspace{0.8em}
\noindent\textbf{\large 二、以衰償損}
\vspace{0.5em}

\noindent 以衰償損者，按各人之過失大小列衰分配賠償也。各以其過失之比例為列衰，
過失大者償多，過失小者償少。

\begin{problem}[二]
今有牛、馬、羊食人苗。苗主責之粟五斗。羊主曰：「我羊食半馬。」馬主曰：「我馬食半牛。」今欲衰償之，問各出幾何？
\end{problem}

\begin{answer}
牛主出二斗八升、七分升之四。馬主出一斗四升、七分升之二。羊主出七升、七分升之一。
\end{answer}

\begin{method}
術曰：置牛四、馬二、羊一，各自為列衰，副并為法。以五斗乘未并者各自為實。實如法得一斗。
\end{method}

\begin{liuhuicomm}
羊食半馬，馬食半牛，則牛食四分，馬食二分，羊食一分，是為列衰。副并得七為法。以五斗乘未并者，牛得二十，馬得一十，羊得五，各自為實。實如法而一，即各出之數也。

按：此題列衰之推求，尤見巧思。羊食半馬，則羊之過為馬之半；馬食半牛，則馬之過為牛之半。以牛之過為四分，則馬為二分，羊為一分。以此為列衰，副并得七為法。以五斗乘未并者，牛四乘五得二十，馬二乘五得一十，羊一乘五得五，各為實。以法七除各實：牛二十除以七得二又七分之六，即二斗八升又七分升之四（因七分之六斗即八升又七分升之四）；馬十除以七得一又七分之三，即一斗四升又七分升之二（因七分之三斗即四升又七分升之二）；羊五除以七得七分之五，即七升又七分升之一（因七分之五斗即七升又七分升之一）。此以比例推列衰之妙也。
\end{liuhuicomm}

\vspace{0.8em}
\noindent\textbf{\large 三、以資產分稅}
\vspace{0.5em}

\noindent 以資產分稅者，按各人之資產多少列衰分配賦稅也。資產多者稅多，資產少者稅少，
此均平之術也。

\begin{problem}[三]
今有甲持錢五百六十，乙持錢三百五十，丙持錢一百八十，凡三人俱出關，關稅百錢。欲以錢數多少衰出之，問各幾何？
\end{problem}

\begin{answer}
甲出五十一錢、一百九分錢之四十一。乙出三十二錢、一百九分錢之一十二。丙出一十六錢、一百九分錢之五十六。
\end{answer}

\begin{method}
術曰：各置錢數為列衰，副并為法，以百錢乘未并者，各自為實，實如法得一錢。
\end{method}

\begin{liuhuicomm}
甲五百六十，乙三百五十，丙一百八十，各以其錢數為列衰。副并得一千九十為法。以百錢乘未并者，甲得五萬六千，乙得三萬五千，丙得一萬八千，各自為實。實如法而一，即各出稅錢也。

按：此術以各人之資產數為列衰，資產多者列衰大，資產少者列衰小。
以此分配關稅，使富者多納，貧者少納，此均平之道也。
甲列衰五百六十，乙三百五十，丙一百八十，副并得一千九十為法。
以所分百錢乘未并者：甲一百乘五百六十得五萬六千，
乙一百乘三百五十得三萬五千，丙一百乘一百八十得一萬八千，各為實。
以法一千九十除各實：甲五萬六千除以一千九十得五十一又一千零九分之四十一，
即五十一錢一百九分錢之四十一（因一千零九約去公因數九，得一百零九）；
乙三萬五千除以一千九十得三十二又一千零九分之一十二，
即三十二錢一百九分錢之一十二；
丙一萬八千除以一千九十得一十六又一千零九分之五十六，
即一十六錢一百九分錢之五十六。
\end{liuhuicomm}

\vspace{0.8em}
\noindent\textbf{\large 四、等比衰分}
\vspace{0.5em}

\noindent 等比衰分者，以等比數列為列衰也。如倍數遞增、半數遞減之類，
皆以等比為列衰而分配之。

\begin{problem}[四]
今有女子善織，日自倍，五日織五尺。問日織幾何？
\end{problem}

\begin{answer}
初日織一寸、三十一分寸之十九。次日織三寸、三十一分寸之七。次日織六寸、三十一分寸之十四。次日織一尺二寸、三十一分寸之二十八。次日織二尺五寸、三十一分寸之二十五。
\end{answer}

\begin{method}
術曰：置一、二、四、八、十六為列衰，副并為法，以五尺乘未并者，各自為實，實如法得一尺。
\end{method}

\begin{liuhuicomm}
日自倍者，初日一，次日二，第三日四，第四日八，第五日十六，是為列衰。副并得三十一為法。以五尺乘未并者，初日得五，次日得一十，第三日得二十，第四日得四十，第五日得八十，各自為實。實如法而一，即各日織之數也。

按：此題以等比數列為列衰，尤見衰分之變化。女子善織，日自倍，即每日織布量為前日之二倍。以初日為一，次日為二，第三日為四，第四日為八，第五日為十六，此以二為公比之等比數列也。副并得三十一為法。以五尺即五十寸乘未并者：初日一乘五十得五十，次日二乘五十得一百，第三日四乘五十得二百，第四日八乘五十得四百，第五日十六乘五十得八百，各為實。以法三十一除各實：初日五十除以三十一得一又三十一分之十九，即一寸三十一分寸之十九；次日一百除以三十一得三又三十一分之七，即三寸三十一分寸之七；第三日二百除以三十一得六又三十一分之十四，即六寸三十一分寸之十四；第四日四百除以三十一得十二又三十一分之二十八，即一尺二寸三十一分寸之二十八；第五日八百除以三十一得二十五又三十一分之二十五，即二尺五寸三十一分寸之二十五。以此驗之，五數相加恰為五尺，術無誤也。
\end{liuhuicomm}

\vspace{0.8em}
\noindent\textbf{\large 五、以算數發傜}
\vspace{0.5em}

\noindent 以算數發傜者，按各鄉之戶口算數列衰分配徭役也。算數多者役多，
算數少者役少，此均役之法也。

\begin{problem}[五]
今有北鄉算八千七百五十八，西鄉算七千二百三十六，南鄉算八千三百五十六，凡三鄉，發傜三百七十八人。欲以算數多少衰出之，問各幾何？
\end{problem}

\begin{answer}
北鄉遣一百三十五人、一萬二千一百七十五分人之一萬一千六百三十七。西鄉遣一百一十二人、一萬二千一百七十五分人之四千四。南鄉遣一百二十九人、一萬二千一百七十五分之八千七百九。
\end{answer}

\begin{method}
術曰：各置算數為列衰，副并為法，以所發傜人數乘未并者，各自為實，實如法得一人。
\end{method}

\begin{liuhuicomm}
按：此術以各鄉之算數為列衰。北鄉算八千七百五十八，西鄉算七千二百三十六，
南鄉算八千三百五十六，副并得二萬四千三百五十為法。
以所發傜人數三百七十八乘未并者：
北鄉三百七十八乘八千七百五十八，約為三百七十八乘八千七百五十八，
得三千三百一十萬零五千二百二十四，為實。以法二萬四千三百五十除之，
得一百三十五又二萬四千三百五十分之一萬一千六百三十七，
約之即一萬二千一百七十五分之一萬一千六百三十七（約去公因數二）。
以此法算西鄉、南鄉亦然。
\end{liuhuicomm}

\begin{problem}[六]
今有稟粟，大夫、不更、簪裹、上造、公士，凡五人，一十五斗。今有大夫一人後來，亦當稟五斗。倉無粟，欲以衰出之，問各幾何？
\end{problem}

\begin{answer}
大夫出一斗、四分斗之一。不更出一斗。簪裹出四分斗之三。上造出四分斗之二。公士出四分斗之一。
\end{answer}

\begin{method}
術曰：各置所稟粟斛斗數，爵次均之，以為列衰，副并而加後來大夫亦五斗，得二十以為法。以五斗乘未并者各自為實。實如法得一斗。
\end{method}

\begin{liuhuicomm}
五人各五斗，合稟二十五斗。今有十五斗，不足十斗。後來大夫亦當稟五斗，并之得二十斗為法。以所不足五斗乘未并者，大夫得二十五，不更得二十，簪裹得一十五，上造得一十，公士得五，各自為實。實如法而一，即各出之數也。

按：此題稍變其術。五人當稟二十五斗，今只有十五斗，不足十斗。
後來大夫亦當稟五斗，此時倉中無粟，須令五人各出其所應得者。
以各人所稟之數為列衰：大夫五、不更四、簪裹三、上造二、公士一，
副并得十五，加後來大夫之五，得二十為法。以所不足五斗乘未并者：
大夫五乘五得二十五，不更五乘四得二十，簪裹五乘三得十五，
上造五乘二得十，公士五乘一得五，各為實。以法二十除各實：
大夫二十五除以二十得一又四分之一，即一斗四分斗之一；
不更二十除以二十得一斗；簪裹十五除以二十得四分之三斗；
上造十除以二十得四分之二斗；公士五除以二十得四分之斗。
合之得五斗，恰為後來大夫所應稟之數。
\end{liuhuicomm}

\begin{problem}[七]
今有稟粟五斛，五人分之，欲令三人得三，二人得二。問各幾何？
\end{problem}

\begin{answer}
三人，人得一斛一斗五升、十三分升之五。二人，人得七斗六升、十三分升之十二。
\end{answer}

\begin{method}
術曰：置三人，人三；二人，人二，為列衰。副并為法。以五斛乘未并者，各自為實。實如法得一斛。
\end{method}

\begin{liuhuicomm}
按：此題列衰之置，又為一變。三人各得三，二人各得二。
故列衰為三、三、三、二、二。副并得十三為法。
以五斛即五十升乘未并者：三人各三乘五十得一百五十，
二人各二乘五十得一百，各為實。以法十三除各實：
一百五十除以十三得一十一又十三分之七，即一斛一斗五升十三分升之五；
一百除以十三得七又十三分之九，即七斗六升十三分升之十二。
此以不同之數為列衰也。
\end{liuhuicomm}

\vspace{0.8em}
\noindent\textbf{\large 六、返衰術}
\vspace{0.5em}

\noindent 返衰術者，與正衰相反之分配術也。正衰以級高者分多，返衰以級高者分少。
其術列置衰而令相乘，動者為不動者衰。返衰之用，在於貴者出少、賤者出多之時。

\begin{method}
返衰術曰：列置衰而令相乘，動者為不動者衰。
\end{method}

\begin{liuhuicomm}
返衰者，反其列衰也。令相乘者，以多乘少，以少乘多，各為其率。動者為不動者衰，言各以不動者為法，動者為實也。

按：返衰術者，衰分術之變也。正衰以列衰之數為率，級高者衰大，所得多。
返衰則反是，級高者衰小，所得少。其術列置衰而令相乘者，
以原列衰各數相乘為新列衰也。如正衰為五、四、三、二、一，
則返衰為一乘二乘三乘四、一乘二乘三乘五、一乘二乘四乘五、
一乘三乘四乘五、二乘三乘四乘五，即二十四、三十、四十、六十、一百二十。
此返衰之新列衰也。以此新列衰分配，則級高者所得少，級低者所得多，
與正衰相反矣。
\end{liuhuicomm}

\begin{problem}[八]
今有大夫、不更、簪褭、上造、公士，凡五人，共出百錢。欲令高爵出少，以次漸多，問各幾何？
\end{problem}

\begin{answer}
大夫出八錢、一百三十七分錢之一百四。不更出一十錢、一百三十七分錢之一百三十。簪褭出一十四錢、一百三十七分錢之八十二。上造出二十一錢、一百三十七分錢之一百二十三。公士出四十三錢、一百三十七分錢之一百九。
\end{answer}

\begin{method}
術曰：置爵數各自為衰，而返衰之，副并為法。以百錢乘未并者各自為實。實如法得一錢。
\end{method}

\begin{liuhuicomm}
高爵出少，以次漸多，故令大夫五、不更四、簪裹三、上造二、公士一，返衰之。大夫得一百二十，不更得六十，簪裹得四十，上造得三十，公士得二十四，是為列衰。副并得二百七十四為法。以百錢乘未并者，各自為實。實如法而一，即各出錢數也。

按：此題返衰之置，以爵數五、四、三、二、一為正衰，返衰之則以各數相乘。
大夫返衰：一乘二乘三乘四得二十四（以大夫五為不動，動一、二、三、四相乘）；
不更返衰：一乘二乘三乘五得三十；
簪裹返衰：一乘二乘四乘五得四十；
上造返衰：一乘三乘四乘五得六十；
公士返衰：二乘三乘四乘五得一百二十。
以此為列衰：二十四、三十、四十、六十、一百二十。
副并得二百七十四為法。以百錢乘未并者：大夫百乘二十四得二千四百，
不更百乘三十得三千，簪裹百乘四十得四千，上造百乘六十得六千，
公士百乘一百二十得一万二千，各為實。以法二百七十四除各實：
大夫二千四百除以二百七十四得八又二百七十四分之一百四十八，
約之即八錢一百三十七分錢之一百四（約去公因數二）；
不更三千除以二百七十四得一十又二百七十四分之一百三十，
約之即一十錢一百三十七分錢之一百三十（約去公因數二）；
以此類推，得各出之數。此返衰之妙也。
\end{liuhuicomm}

\begin{problem}[九]
今有甲持粟三升，乙持糲米三升，丙持糲飯三升。欲令合而分之，問各幾何？
\end{problem}

\begin{answer}
甲二升、一十分升之七。乙四升、一十分升之五。丙一升、一十分升之八。
\end{answer}

\begin{method}
術曰：以粟率五十、糲米率三十、糲飯率七十五為衰、而返衰之，副并為法。以九升乘未并者各自為實。實如法得一升。
\end{method}

\begin{liuhuicomm}
甲持粟率五十，乙持糲米率三十，丙持糲飯率七十五，返衰之。甲得四千五百，乙得七千五百，丙得三千，是為列衰。副并得一萬五千為法。以九升乘未并者，各自為實。實如法而一，即各得之數也。

按：此題以各物之率為正衰，返衰之則以各率之乘積為新衰。
粟率五十、糲米率三十、糲飯率七十五，返衰之：
甲粟返衰：三十乘七十五得二千二百五十；
乙糲米返衰：五十乘七十五得三千七百五十；
丙糲飯返衰：五十乘三十得一千五百。
以此為列衰：二千二百五十、三千七百五十、一千五百。
副并得七千五百為法（可約之）。以九升乘未并者各為實，
實如法而一，即各得之數。劉徽注中所云甲得四千五百、乙得七千五百、
丙得三千，即將上述列衰各乘二所得，法亦乘二，故其實同也。
\end{liuhuicomm}

\vspace{0.8em}
\noindent\textbf{\large 七、均輸雜題}
\vspace{0.5em}

\noindent 衰分之章，不獨論比例分配而已，亦雜有均輸、利息、雇傭之題，
皆以今有術、衰分術為基礎而變化之。

\begin{problem}[一０]
今有絲一斤，價直二百四十。今有錢一千三百二十八，問得絲幾何？
\end{problem}

\begin{answer}
五斤八兩一十二銖、三分銖之四。
\end{answer}

\begin{method}
術曰：以一斤價數為法，以一斤乘今有錢數為實，實如法得絲數。
\end{method}

\begin{liuhuicomm}
按：此術猶今有之義。以一斤價數為法，以所求率乘所有數為實，實如法而一也。
一斤價二百四十為法，一斤乘今有錢一千三百二十八得一千三百二十八，為實。
實如法而一，得五斤餘一百二十八。以一斤十六兩乘餘數，以二百四十除之，
得八兩餘一百二十八。以一两二十四銖乘餘數，以二百四十除之，
得一十二銖餘一百九十二。以三分銖之四命之，即五斤八兩一十二銖三分銖之四也。
\end{liuhuicomm}

\begin{problem}[一一]
今有絲一斤價直三百四十三。今有絲七兩一十二銖，問得錢幾何？
\end{problem}

\begin{answer}
一百六十一錢、三十二分錢之二十三。
\end{answer}

\begin{method}
術曰：以一斤銖數為法，以一斤價數，乘七兩一十二銖為實。實如法得錢數。
\end{method}

\begin{liuhuicomm}
按：一斤三百八十四銖（一斤十六兩，一兩二十四銖，十六乘二十四得三百八十四）。
以三百八十四為法。七兩一十二銖，七兩即一百六十八銖，加一十二銖得一百八十銖。
以一斤價三百四十三乘一百八十得六万一千七百四十，為實。
實如法而一，以三百八十四除之，得一百六十一餘二百二十三，
即一百六十一錢又三百八十四分錢之二百二十三，約之即三十二分錢之二十三也。
\end{liuhuicomm}

\begin{problem}[一二]
今有縑一丈價直一百二十六。今有縑一匹九尺五寸，問得錢幾何？
\end{problem}

\begin{answer}
六百三十三錢、五分錢之三。
\end{answer}

\begin{method}
術曰：以一丈寸數為法，以價錢數乘今有縑寸數為實，實如法得錢數。
\end{method}

\begin{problem}[一三]
今有布一匹，價直一百二十三。今有布二丈七尺，問得錢幾何？
\end{problem}

\begin{answer}
八十四錢、分錢之三。
\end{answer}

\begin{method}
術曰：以一匹尺數為法，今有布尺數乘價錢為實，實如法得錢數。
\end{method}

\begin{problem}[一四]
今有素一匹一丈，價直六百二十五。今有錢五百，問得素幾何？
\end{problem}

\begin{answer}
得素四丈。
\end{answer}

\begin{method}
術曰：以價直為法，以一匹一丈尺數乘今有錢數為實。實如法得素數。
\end{method}

\begin{liuhuicomm}
此術猶今有之義。以價直為法，以所求率乘所有數為實，實如法而一也。
按：素一匹一丈，即四丈也（一匹四丈）。價直六百二十五為法。
四丈即四十尺，乘今有錢五百得二萬，為實。實如法而一，
以六百二十五除二萬得三十二尺，即四丈也。
\end{liuhuicomm}

\begin{problem}[一五]
今有與人絲一十四斤，約得縑一十斤。今與人絲四十五斤八兩，問得縑幾何？
\end{problem}

\begin{answer}
三十二斤八兩。
\end{answer}

\begin{method}
術曰：以一十四斤兩數為法，以一十斤乘今有絲兩數為實，實如法得縑數。
\end{method}

\begin{liuhuicomm}
按：此術亦今有術也。以一十四斤兩數為法，一十四斤即二百二十四兩（十四乘十六）。
以一十斤乘今有絲兩數為實。今有絲四十五斤八兩，即七百二十八兩（四十五乘十六加八）。
以十斤即一百六十兩乘七百二十八得十一萬六千四百八十，為實。
實如法而一，以二百二十四除之，得五百二十兩，即三十二斤八兩（三十二乘十六得五百一十二，
加八得五百二十）。
\end{liuhuicomm}

\begin{problem}[一六]
今有絲一斤，耗七兩。今有絲二十三斤五兩，問耗幾何？
\end{problem}

\begin{answer}
一百六十三兩四銖半。
\end{answer}

\begin{method}
術曰：以一斤展十六兩為法，以七兩乘今有絲兩數為實，實如法得耗數。
\end{method}

\begin{liuhuicomm}
按：此術亦比例也。一斤十六兩，耗七兩，即每十六兩耗七兩也。
以十六兩為法，七兩乘今有絲兩數為實。今有絲二十三斤五兩，
即三百七十三兩（二十三乘十六加五）。以七兩乘三百七十三得二千六百一十一，為實。
實如法而一，以十六除之，得一百六十三兩餘三兩。三兩即七十二銖（三乘二十四），
以十六除之得四銖半，即一百六十三兩四銖半也。
\end{liuhuicomm}

\begin{problem}[一七]
今有生絲三十斤，乾之，耗三斤十二兩。今有乾絲一十二斤，問生絲幾何？
\end{problem}

\begin{answer}
一十三斤一十一兩十銖、七分銖之二。
\end{answer}

\begin{method}
術曰：置生絲兩數，除耗數，餘，以為法。三十斤乘乾絲兩數為實。實如法得生絲數。
\end{method}

\begin{liuhuicomm}
此術猶今有之義。以生絲兩數除耗數，餘為率。三十斤乘乾絲兩數為實，實如法而一，即得生絲之數也。
按：生絲三十斤即四百八十兩，耗三斤十二兩即六十兩，除耗餘四百二十兩為法。
三十斤即四百八十兩乘乾絲十二斤即一百九十二兩，得九萬二千一百六十，為實。
實如法而一，以四百二十除之，得二百一十九兩餘一百八十兩。
二百一十九兩即一十三斤一十一兩（十三乘十六得二百零八，加十一得二百一十九）。
餘一百八十兩以四百二十為分母，約之得七分之三。三兩即七十二銖，
以七除之得十銖餘二，即十銖七分銖之二也。
\end{liuhuicomm}

\begin{problem}[一八]
今有田一畝，收粟六升、太半升。今有田一頃二十六畝一百五十九步，問收粟幾何？
\end{problem}

\begin{answer}
八斛四斗四升、一十二分升之五。
\end{answer}

\begin{method}
術曰：以畝二百四十步為法，以六升、太半升乘今有田積步為實，實如法得粟數。
\end{method}

\begin{liuhuicomm}
按：此術以畝產量為率，以田積步為所有數，求所收粟數。
畝法二百四十步為法。六升太半升即六升又二分之一升，即二分之十三升。
以二分之十三升乘田積步為實。田一頃二十六畝一百五十九步，
一頃百畝，共一百二十六畝。一百二十六畝乘二百四十步得三萬零二百四十步，
加一百五十九步得三萬零三百九十九步，為田積步。
以二分之十三乘三萬零三百九十九，約得十九萬七千五百九十四分之一，
為實。實如法而一，以二百四十除之，得八斛四斗四升一十二分升之五也。
\end{liuhuicomm}

\begin{problem}[一九]
今有取保一歲，價錢二千五百。今先取一千二百，問當作日幾何？
\end{problem}

\begin{answer}
一百六十九日、二十五分日之二十三。
\end{answer}

\begin{method}
術曰：以價錢為法，以一歲三百五十四日乘先取錢數為實，實如法得日數。
\end{method}

\begin{liuhuicomm}
按：此術亦今有術也。取保一歲價錢二千五百為法，一歲三百五十四日為所有率，
先取錢一千二百為所有數。以所有率乘所有數為實，實如法而一，
即當作之日數。三百五十四乘一千二百得四十二萬四千八百，為實。
以二千五百為法，實如法而一，得一百六十九日餘二千三百，
即一百六十九日二千五百分日之二千三百，約之即二十五分日之二十三也。
\end{liuhuicomm}

\begin{problem}[二０]
今有貸人千錢，月息三十。今有貸人七百五十錢，九日歸之，問息幾何？
\end{problem}

\begin{answer}
六錢、四分錢之三。
\end{answer}

\begin{method}
術曰：以月三十日，乘千錢為法。以息三十乘今所貸錢數，又以九日乘之，為實。實如法得一錢。
\end{method}

\begin{liuhuicomm}
此術以月三十日乘千錢為法者，得一萬，是為貸錢一萬，一日息三十也。以息三十乘今所貸七百五十錢，得二萬二千五百，又以九日乘之，得二十萬二千五百，為實。實如法得一錢，即九日之息也。

按：此術亦今有術之變。月三十日乘千錢得三萬為法。以息三十乘今所貸七百五十得二萬二千五百，又以九日乘之得二十萬二千五百為實。實如法而一，以三萬除二十萬二千五百，得六錢餘二千五百，即六錢三萬分錢之二千五百，約之即四分錢之三也。

此題計息之法，以日數乘錢數，以月率除之，即得利息。此為中國古代利息計算之最早記載，開後世金融數學之先河。
\end{liuhuicomm}

\vspace{0.8em}
\noindent\textbf{\large 八、衰分章小結}
\vspace{0.5em}

\begin{quoting}
衰分之章，凡二十問。論比例分配之法，以列衰為綱，以衰分術為領。
衰分術曰：各置列衰，副并為法，以所分乘未并者各自為實，實如法而一。
此術與今有術相通，而變化尤多。

本章所列之題，或以爵次分物，或以衰償損，或以資產分稅，或以算數發傜，
或以均輸雜算，皆衰分術之應用也。又有返衰術，以反其列衰，
使高爵出少、以次漸多，與正衰相反，而術理則一。

劉徽之注，闡明衰分之理，謂「衰分者，差分也。列衰者，列其差數以為率也」，
此衰分術之精髓。李淳風等又按密率校之，使精度益進。
衰分之術，於今之分配問題、比例計算、稅賦分攤，無不可用，其義博矣。

方田、粟米、衰分三卷，為《九章算術》之基礎。方田論面積與分數，
粟米論比例與交換，衰分論分配與均輸。三卷相為表里，奠定中國古代數學之基礎。
劉徽之注，發微闡幽，多所發明，尤以割圓術為最，開極限思想之先河，
誠千古不朽之作也。
\end{quoting}

\vspace{1em}
\begin{center}
\rule{0.5\textwidth}{0.4pt}
\end{center}

\vspace{1em}

\begin{center}
\textit{（以上為《九章算術》卷第一至卷第三之全文，含劉徽注及李淳風等注釋。卷一方田論田地面積之計算及分數運算；卷二粟米論谷物交換之比例計算；卷三衰分論按比率分配之法。此三卷為《九章算術》之基礎，示中國古代數學之精髓。）}
\end{center}

\end{document}
